\section{Method}
\label{sec:method}
\subsection{Reconstruction of the LV cavity}
\label{sec:lv-method}
In the test segmentations, the LV cavity is background enclosed by the myocardium. A natural reconstruction is to take, slice by slice, the holes of the myocardium mask (\emph{fill}). This fails on slices where the myocardial ring is open, typically the most basal one, where the myocardium does not fully surround the cavity: the slice then gets no cavity at all. We add a fallback (\emph{fill+hull}): on a slice with myocardium but no hole, the cavity is the convex hull of the myocardium minus the myocardium itself, restricted to background voxels. On apical slices, where the myocardium is a solid cap, this hull is a thin sliver; we keep it only if its area is at least 25\% of the myocardium area. Both methods are evaluated on the training set, where the true cavity is known, by removing it and reconstructing it (Section~\ref{sec:lv}).

\subsection{Features}
Table~\ref{tab:features} lists the features. The 21 \emph{volumetric} features are classical biomarkers: LV and RV volumes at ED and ES, ejection fractions, the myocardial mass (myocardium volume times 1.05 g/mL), volumes and mass indexed by the body surface area (BSA) computed with the Mosteller formula \citep{mosteller1987}, RV / LV ratios, height and weight.

None of them describes the shape of the myocardial wall, although wall thickness defines HCM \citep{elliott2014} and an infarct leaves a thin, akinetic segment \citep{thygesen2019}. We therefore add 11 \emph{thickness} features, in the spirit of the ACDC entries \citep{isensee2017,wolterink2017,khened2017}. On each slice that contains both myocardium and cavity, we compute the Euclidean distance transform (in mm) of the region formed by the myocardium and the cavity, and read it at every myocardium pixel that touches the cavity: this is the distance from the inner border of the wall to its outer border, i.e.\ the local wall thickness. For ED and ES, we summarize these samples by their mean and maximum, the largest per-slice mean, the mean of the per-slice standard deviations (how uneven the wall is within a slice) and the standard deviation of the per-slice means (how uneven it is along the axis). The systolic thickening is the ratio of the mean thicknesses at ES and ED. We compare three feature sets: the 21 volumetric features, the same plus the 11 thickness features, and a \emph{compact} set of 12 BSA-indexed and thickness features without raw height and weight.

\begin{table}[!htb]
\centering
\small
% Written by hand from final/cardiac/features.py (definitions in Appendix, tabs/supp/features_full.tex).
\begin{tabular}{p{0.2\linewidth}rp{0.62\linewidth}}
\toprule
Group & Count & Features\\
\midrule
LV & 6 & ED and ES volumes, ejection fraction, stroke volume, ED and ES volumes / BSA\\
RV & 5 & ED and ES volumes, ejection fraction, ED and ES volumes / BSA\\
Myocardium & 5 & mass (mean of ED and ES), mass / BSA, ES / ED mass, ED mass / LV ED volume, mean mass / LV ED volume\\
RV / LV & 2 & volume ratios at ED and at ES\\
Patient & 3 & height, weight, body surface area (Mosteller)\\
\midrule
Wall thickness & 11 & at ED and at ES: mean, maximum, largest per-slice mean, mean of per-slice standard deviations, standard deviation of per-slice means; systolic thickening (ES / ED mean thickness)\\
\bottomrule
\end{tabular}
\caption{Features. The first 21 are the \emph{volumetric} features. The \emph{compact} set keeps 12 of them: the 6 indexed volumes and ejection fractions, mass / BSA, the RV / LV volume ratio at ED, mean mass / LV ED volume, the largest per-slice mean and the mean per-slice standard deviation of the ED wall thickness, and the systolic thickening.}
\label{tab:features}

\end{table}

\subsection{Classifiers}
Every model is a pipeline that standardizes the features on the training data. We compare logistic regression, linear and RBF SVMs, $k$ nearest neighbors, random forests, extremely randomized trees and XGBoost, with the hyper-parameter grids of Appendix~\ref{app:bench}, selected by the inner cross-validation.

\paragraph{MLP ensemble}
We use an ensemble of 50 MLPs with 4 hidden layers of 32 units, batch normalization and leaky ReLUs, each trained on a random 75\% of the training patients and a random two thirds of the features, for 200 epochs with Adam (learning rate $10^{-3}$, batches of 30). The ensemble averages their class probabilities. We train the 50 networks as one batched network: the weights of each layer are stacked into a tensor, each network sees its own patients and batch-normalization statistics, and dropped features are set to zero, which is equivalent to removing them. This gives the same result as 50 separate networks about 50 times faster, which makes nested cross-validation affordable.

\paragraph{Soft vote}
The MLP ensemble is combined with a random forest or a linear SVM by averaging their probabilities with a weight $\alpha$ on the MLP:
\begin{equation}
  p(y\mid x) = \alpha\, p_{\mathrm{MLP}}(y\mid x) + (1-\alpha)\, p_{\mathrm{other}}(y\mid x).
  \label{eq:softvote}
\end{equation}
We choose $\alpha\in\{0, 0.1, \dots, 1\}$ inside each training fold, from the out-of-fold probabilities of an inner 5-fold cross-validation, breaking ties towards $\alpha=0.5$; the two models are then refitted on the whole training fold. The SVM probabilities come from Platt scaling \citep{platt1999}.
